\scanpage{11}\mspage{10}
En \deleted{\unsure{l}}effet, c'est comme pour les sous-familles admissibles d'une famille (ici $\Phi\cup\deleted{\unsure{\Phi}}\Psi$) -- cf. Scholie page 7.

\margin{\uline{Cor} 3 Si $\Phi,\Psi$ comp., alors pour $\Phi',\deleted{\unsure{\Psi}}\leq\Phi,\ \Psi'\leq\Psi$ (parties fermées) $\Phi'$ et $\Psi'$, \unsure{st} admissibles, et $\Phi'\cup\Psi'$ est fermé (ss-famille adm.) dans $\Phi\cup\Psi$.}

% slanted
\margin{On appellera \unsure{ainsi|aussi} les éléments de $\unsure{\mathcal F|\mathscr F}$ les \uline{``figures \unsure{ensemblistes}''} de $L$}

Soit $\mathcal F$ l'ens des \deleted{familles} \unsure{des|les} parties \uline{admissibles} $\Phi$ de $\unsure{\mathcal P|\mathfrak P}(L)$. On ordonne $\mathcal F$ par
\[
\begin{aligned}
\Phi\leq\Psi\quad\overset{\mathrm{def}}{\Longleftrightarrow}\quad &\Phi\ \text{est une ss-}\\
&\text{famille admissible}\\
&\text{de }\Psi.
\end{aligned}
\]

Alors on a

\uline{C 1} $\mathcal F$ admet des Sup majorés, car si $\Phi_i\in\mathcal F,\ \Phi_i\leq\deleted{\Phi}\Psi$ des ss-ens. fermés d'un $=$ $\deleted{\Phi}\Psi$, alors $\bigcup_i\Phi_i$ est un sous-ens. fermé de $\Psi$, qui est \unsure{le|la} la sup. C'est immédiat par \unsure{le sorite|les sorites} sur les sous-familles \unsure{équiv.} (\unsure{de ses|i.e.|les|deux} parties fermées) d'un $\Psi$

% slanted
\margin{La Sup pour la relation $\leq$ (qui est ici $\subset$) est $\bigcup_i$}

\deleted{\unsure{Oui|Or|O.K.}, un peu \unsure{vif|cf.}, il faut \unsure{passer|préciser} par des Sup finis, en notant que $\emptyset\subset\mathcal P(L)$ est admissible (! cas du Sup \inserted{de la \unsure{fam.}} vide !) -- la seule famille admissible $\underset{\in\mathcal F}{\Phi}$ telle que $\lvert\Phi\rvert=\emptyset$ et que si $\Phi,\Psi,\Sigma$ sont tels que}

% slanted
\margin{\unsure{explicitons|expliciter} \unsure{cf.} \uuline{Prop} ou Cor}

% This line has a horizontal cancellation in addition to the diagonal strokes across the draft.
\deleted{$\Phi$ et $\Psi$) est admissible, alors $\Sigma$ \unsure{ai|est}}

\deleted{($\Phi,\Psi$) compatibles, alors $\Sigma$ est compatible avec $\Phi\unsure{\cup|et}\Psi$ \unsure{sss|ssi} il l'est avec $\Phi$ et avec $\Psi$}
